\section{Solving a Unique Games problem with a partition function}\label{sec:solveCUG}
%
In this section we define a partition function $Z(G;w)$ and describe how to solve CUG instances via knowledge of the partition function, leading to proofs of Theorems~\ref{thm:main,interp} and~\ref{thm:main,cluster}. 
We also observe how CUG naturally arises from UG in this context. 
%
A partition function is a mathematical object that encodes as a polynomial some weighted substructures in a graph. 
The general definition arises in the statistical physics of spin systems, and important examples include the independence polynomial and matching polynomials of a graph, see e.g.~\cite{Barvinok16}. 
Here we will only describe the partition function we define to study Unique Games instances, which is closely related to the Potts and random cluster models.
%
\begin{definition}\label{def:Z}
Given a Unique Games instance $G=(V, E, \pi)$, the partition function $Z(G;w)$ is a polynomial in a parameter $w\in\CC$ given as a sum of terms $w^{i}$ for each coloring of the graph with the alphabet $[k]$ that has $i$ satisfied constraints (i.e.\ that has value $i/|E|$).
That is,
\begin{equation}
Z(G; w) := \sum_{\{x_u\}_{u\in V}\in[k]^V}\;\prod_{\substack{(u,v)\in E,\\ x_v=\pi_{uv}(x_u)}}w\,,
\end{equation}
where the sum us over all assignments of colors in $[k]$ to the labels $\{x_u\}_{u\in V}$.
\end{definition}
%
If the permutations $\pi_{uv}$ are all the identity, then a satisfied constraint corresponds to a \emph{monochromatic edge}: both endpoints of the edge received the same color. 
In this case the above $Z(G;w)$ corresponds to the partition function of the Potts model from statistical physics. 
The relevance to Unique Games problems arises from the fact that the cases of the promise in (C)UG give contrasting bounds on $Z(G;w)$. 
When $w\ge 1$ is real, in the case that a highly satisfying assignment is guaranteed to exist we have a lower bound, and in the case that no highly satisfying assignments exist we have an upper bound. 
When the upper bound is less than the lower bound, at most one of the bounds can hold for any given instance, so knowledge of $Z(G;w)$ immediately solves the problem. 
If the bounds are sufficiently far apart an approximate value of $Z(G;w)$ suffices.  
%
\begin{lemma}\label{lem:solveCUG}
Consider an instance $G=(V,E)$ of \CUG[f,k,\eps,\delta], and let $\alpha>0$. 
Then to solve the instance it suffices to know any value $\xi$ satisfying $e^{-\alpha} \le Z(G;w)/\xi \le e^\alpha$ for any real $w$ such that
\[
\log w > \frac{|V|\log(1/f)+2\alpha}{(1-\eps-\delta)|E|}\,.
\]
In the case that $G$ has average degree $\Delta$ (so $2|E|=\Delta|V|$), and $\alpha=C|V|$ this becomes 
\[
\log w > \frac{2}{1-\eps-\delta}\frac{\log(1/f)+2C}{\Delta}\,.
\]
\end{lemma}
%
\begin{proof}
Consider any real $w\ge 1$.
Then if there are $(fk)^{|V|}$ colorings of $G$ with value $1-\eps$ we have 
\[
e^{\alpha}\xi \ge Z(G;w) \ge (fk)^{|V|}w^{(1-\eps)|E|}\,,
\]
and if every coloring has value at most $\delta$ we have 
\[
e^{-\alpha}\xi \le Z(G;w)\le k^{|V|}w^{\delta|E|}\,.
\] 
Since these bounds go in opposite directions, knowledge of $\xi$ immediately yields an solution to the problem when the implied intervals for $\xi$ are disjoint. 
This occurs precisely for $w$ as in the statement of the lemma. 
\end{proof}
%
We now see that increasing $f$ allows CUG instances to be solved via $Z(G;w)$ for smaller $w$, and this is how the Count variant of Unique Games naturally arises. 
We are unable to approximate $Z(G;w)$ with $w$ large enough to permit $f=1/k$ (i.e.\ refute UGC), but by slightly increasing $f$ we bring $w$ into a range amenable to our methods. 
In Section~\ref{sec:conc} we discuss the problem of how large $f$ can be while \CUG[f] is still equivalent to UG, and identify natural barriers to using algorithms for larger $w$.
To obtain Theorems~\ref{thm:main,interp} and~\ref{thm:main,cluster} we need a pair of algorithms and some calculations. 
%
\begin{theorem}\label{thm:alg,interp}
Let $k\in \mathbb{N}_{\geq 3}$, $\Delta\in \mathbb{N}_{\geq 3}$, and $w^*=1+(\log k -1)/\Delta$. 
Then there exists a deterministic algorithm which, given $\alpha$ satisfying $0<\alpha<\frac{2n}{e\Delta}e^{O(\log^2 k)}$,
and an $n$-vertex \UG[k] instance $G$ of maximum degree at most $\Delta$, computes a number $\xi$ satisfying $e^{-\alpha} \le Z(G;w^*)/\xi \le e^{\alpha}$ in time bounded by
\[
n\exp\left(e^{O(\log^2 k)} \log\left(\frac{n\Delta}{\alpha}e^{O(\log^2 k)}\right)\log(\Delta\sqrt k)\right)\,.
\]
\end{theorem}
%
\begin{theorem}\label{thm:alg,cluster}
Let $\Delta\in \mathbb{N}$ and let $\zeta=8/\sqrt{\Delta}$.
For $k\geq\exp{((18\Delta+4\Delta\log\Delta)/\zeta)}$ and $w^*=\exp((2-\zeta) \log(k)/\Delta)$ there exists an deterministic algorithm, which given $\alpha>0$ and an $n$-vertex \UG[k] instance $G$ of maximum degree at most $\Delta$, computes $\xi$ satisfying $e^{-\alpha}\le |Z(G;w^*)/\xi| \le e^\alpha$ in time bounded by $kn^{O(1)}(n/\alpha)^{O(\Delta)}$.
\end{theorem}
%
These results are proved in Sections~\ref{sec:interp} and~\ref{sec:cluster} respectively. 
In both cases we define a series for $\log Z(G;w)$, show that it converges, and obtain an additive approximation to it by evaluating a truncation of the series. 
Here we give the calculations that show what CUG problems we can solve with these algorithms. 
%
\begin{proof}[Proof of Theorem~\ref{thm:main,interp}]
Take $w^*$ as in Theorem~\ref{thm:alg,interp}.
As $\Delta\to\infty$, and with approximation error $\alpha=Cn$ for any $C<2e^{O(\log^2k)-1}/\Delta$, by Lemma~\ref{lem:solveCUG} we require 
\[
\log w^* = (1-o(1))\frac{\log k}{\Delta} > \frac{2}{1-\eps-\delta}\frac{\log(1/f)+2C}{\Delta}\,.
\]
For large enough $\Delta$, $C=\log(k)/\Delta$ is valid in Theorem~\ref{thm:alg,interp} and implies the above for any
\[
f \ge k^{\eps+\delta-\frac{1}{2}}\,.\qedhere
\]
\end{proof}
We remark that a very similar calculation gives a result for $k$ growing with $\Delta$, but we present the special case of constant $k$ here as it is instructive of our methods and permits a concise expression for $f$ and the running time.
%
\begin{proof}[Proof of Theorem~\ref{thm:main,cluster}]
Take $w^*$ as in Theorem~\ref{thm:alg,cluster}. 
With $\Delta\ge e^{9/2}$, $\zeta=8/\sqrt\Delta$, $k\ge \Delta^{\Delta^{3/2}}$, and $\alpha=Cn$ for some $C>0$ we choose later, by Lemma~\ref{lem:solveCUG} we require 
\[
\log w = (2-\zeta)\frac{\log k}{\Delta} > \frac{2}{1-\eps-\delta}\frac{\log(1/f)+2C}{\Delta}\,,
\]
which holds when
\[
f > e^{2C}k^{-\frac{1}{2}(2-\zeta)(1-\eps-\delta)}\,.
\]
With $\Delta\ge\Delta_0(\eps,\delta)$ and $C \le \frac{1}{4}(\eps+\delta)\log k$, this holds for $f \ge k^{2\eps+2\delta-1}$.
For large enough $\Delta_0$ (which makes $k$ sufficiently large) we can take e.g. $C=1/2$ and obtain a running time of $kn^{O(1)}e^{O(\Delta)}$.
\end{proof}
%
%
